\section{Introduction}\label{sec:intro}
Fault-tolerant quantum computation requires that physical noise rates be suppressed low enough, and decoding is done fast enough. 
Quantum error correcting codes (QECCs) achieve the former by restricting a larger space of physical 
qubits to a subspace of logical qubits that satisfy some relations.
These relations can be measured without collapsing the encoded quantum state. The design of such codes is just one side of a foundational problem in the path to scalable 
quantum computing. The choice of decoder is another, and the two operate together.
%For binary and Additive White Gaussian Noise (AWGN) channels, low-density parity check (LDPC) codes~\cite{gallager1962low} 
%were adopted in communication and data storage,  
%with a ranging array of decoding .
\noindent Quantum Low Density Parity Check codes (qLDPC)~\cite{panteleev2022asymptotically,bravyi2024high} present an attractive approach as they promise good asymptotic parameters and can be decoded using algorithms such as belief propagation (BP)~\cite{pearl1988probabilistic}.
Furthermore, the implication of \emph{low density} is that the connectivity required between logical qubits is low, an appealing property in some qubit architectures and quantum control hardware.
Practical decoding for fault tolerant quantum computing is likely to constrain the 
decoder design further than qubit connectivity.
Much like the binary setting, the decoder must complete its work under a 
certain time limit, otherwise its decoding latency may kill the entire application~\cite{sella2018fec}.

\noindent A central insight motivating this work is that code design and decoder design are not independent problems, 
especially in decoders that rely on heuristics and approximations of Maximum Likelihood (ML) decoding.
The performance of a given decoder depends strongly on the structure of the code.  
We can ask the question: \textit{given a fixed decoder, which code maximises its performance ?}\\

\noindent A deeper question that motivated this work is: \textit{Assuming deep knowledge of a quantum-device-under-test (qDUT), its error modalities, failure modes etc., how could one use that knowledge to engineer better, stack-integrated, QECC-decoder pair ?}

\noindent While we are still far from answering the latter, we believe this work presents some progress in answering the former. This paper proceeds as follows:\\
Section~\ref{sec:background} is meant to serve as a presentation of the tools we use in this work for a non expert in the corresponding field. We also include a \textbf{dataset} of BB codes we used for supervised learning in a later section.\\
Section~\ref{sec:mdp} describes a \textbf{Markov Decision Process} (MDP) formulation of Bicycle-Bivariate (BB) code design in which the state is a parity check matrix and an action modifies the matrix. Following the action we present a \textbf{reward} function based on the area under a fitted BER physical-error-rate curve, giving a single scalar.\\
\noindent A ``warm start'' strategy, in the form of a predictive model, is presented in Section~\ref{sec:critic}. It is then used in as part of the \textbf{policy architecture} in order to improve sampling efficiency.\\
\noindent Following the related work in Section~\ref{sec:related-work}, we evaluate the training and resulting policy in Section~\ref{sec:evaluation}, and discuss potential further work in Section~\ref{sec:conclusion}.